\chapter[Conclusão]{Conclusão}\label{sec-conclusao}
Este trabalho propôs o uso de transformadas de ondaletas via \emph{lifting scheme} como etapa de pré-processamento para a inferência causal funcional. 
Essa abordagem superou limitações de métodos globais tradicionais, especialmente a dificuldade em capturar comportamentos localmente irregulares e a necessidade de respeitar a causalidade operacional. 
Nossa metodologia nos permitiu ainda estruturar um fluxo de dados para a atualização sequencial e automática das estimativas do \(\mathrm{FATE}\), bem como propor o \(\mathrm{CFATE}\) como medida acumulada do efeito causal ao longo do tempo. 
Ambas as contribuições se valeram de nossa implementação de ondaletas, por meio do pacote \texttt{rLifting} \cite{dasilvaRLiftingHighPerformanceWavelet2026}, desenvolvido para o ecossistema \texttt{R} e disponível no CRAN.

Os resultados das simulações contribuem para evidenciar a superioridade da nossa proposta frente a métodos vigentes na literatura, sobretudo em cenários de alta irregularidade temporal. 
Em experimentos controlados, o uso de ondaletas não apenas reduziu os erros quadráticos, como também garantiu coberturas simultâneas nominais quando aliado ao procedimento de \emph{bootstrap} não paramétrico proposto na Seção \ref{sec-bootstrap-ic}. 
A Simulação 3 mostrou, ainda, que uma malha discreta e não-uniforme de dados não compromete o desempenho do estimador. 
Coberturas e erros com essa grade de dados foram compatíveis com os da grade uniforme da Simulação 1 em todos os graus de irregularidade avaliados. 
Além disso, os resultados obtidos indicam que a perda de performance entre os estimadores \emph{offline} e \emph{online}, quando existente, é mínima, validando a viabilidade do método para aplicações em tempo real.

Ilustramos nosso método aplicando-o aos dados do experimento \emph{Low Carbon London} \cite{schofield2015lcl}.
Com \(N = 5.095\) domicílios e \(T = 17.520\) intervalos de 30 minutos ao longo de 2013, estimamos que a tarifa dinâmica induziu uma redução acumulada de \(-372\) kWh/hh [IC 95\%: \(-534\); \(-209\)].
A análise por grupo ACORN mostrou que grupos mais abastados responderam mais intensamente ao incentivo tarifário, com \emph{Affluent} acumulando \(-486\) kWh/hh e \emph{Adversity} \(-204\) kWh/hh.
Os estimadores \emph{offline} e \emph{online} produziram trajetórias praticamente idênticas, confirmando em dados reais o padrão observado nas simulações.

A despeito desses avanços, o amadurecimento desta linha de pesquisa aponta para questões metodológicas que impõem desafios técnicos e computacionais. 
Algumas questões práticas e escolhas paramétricas foram discutidas na Seção \ref{sec-metodologia-consideracoes}. 
A generalização do estimador \(\mathrm{FATE}\) para cenários de dados mais complexos, no entanto, envolve compromissos de desempenho que merecem maior detalhamento.

Do ponto de vista do \texttt{rLifting}, a frente de evolução mais impactante seria a transição de uma filtragem univariada para uma abordagem multivariada, que nos permitiria incorporar a estrutura de covariância entre o desfecho e variáveis de confusão, com esperado ganho de performance das estimativas estatísticas do \(\mathrm{FATE}\) em troca de maior latência computacional. 
O \texttt{rLifting} já permite um pré-processamento em grades irregulares, mas o estimador do \(\mathrm{FATE}\) ainda exige que essa malha de dados seja compartilhada entre todos os indivíduos. 
Refinar o processo de estimação para relaxar essa suposição aumentaria o escopo do método. 
Somado a isso, a hipótese de linearidade poderia ser relaxada mediante o emprego de métodos de aprendizado de máquina adaptados ao contexto funcional (sob o custo de maior latência computacional). 
A expansão para cenários com covariáveis variantes no tempo e regimes de tratamento dinâmicos aproximaria nossa proposta dos arcabouços discutidos por \citeonline{ying_causality_2024}, também se caracterizaria como um avanço importante, ainda que demande suposições mais fortes para assegurar a identificabilidade. 
Por fim, vale notar que, embora tenhamos seguido a filosofia de modelagem do desfecho de \citeonline{ecker_causal_2024}, seria possível incluir modelos para a atribuição do tratamento, em linha com o proposto por \citeonline{testa_doubly-robust_2025}, conferindo a propriedade de dupla robustez frequentemente buscada em estimadores causais \cite{hernan_causal_2023}.

Qualquer uma dessas frentes serve como proposta para trabalhos futuros na área. 
Alguns desses avanços estariam mais associados à parte de inferência causal funcional, outros à área de processamento de dados; esperançosamente, alguns contribuirão com ambas as frentes, como tentamos fazer ao longo de nosso texto.

Em última análise, este trabalho estabelece uma infraestrutura que reconcilia o rigor da inferência causal funcional com as exigências de desempenho do processamento em tempo real. 
Ao deslocar o eixo do pré-processamento de métodos globais para a localidade das ondaletas via \emph{lifting scheme}, não apenas superamos limitações de métodos tradicionais em cenários de alta irregularidade, mas também viabilizamos uma operabilidade causal até então pouco explorada na literatura causal-funcional. 
Esperamos que as contribuições aqui apresentadas sirvam de alicerce para sistemas de monitoramento contínuo em que a capacidade de discernir relações de causa e efeito, de forma célere e precisa, seja fundamental para a tomada de decisões em ambientes mais dinâmicos.